\section{Empirical convexification in any input dimension}
The objective and training-vector conclusions do not depend on the
one-dimensional Laplace identity. Gaussian entropy coercivity gives
the following extension directly. The fiber and range statements
remain specifically one-dimensional.

\begin{proposition}[A dimension-explicit empirical objective bound]
Let $d\ge1$, $x\in[-1,1]^d$, $\bar x=(x,1)$, and
\[
 z_w(x)=\sum_{j=1}^N v_j(w_j\cdot\bar x)_+,
 \qquad w_j\sim N(0,(\nu/d)I_{d+1}),\qquad
 v_j\sim N(0,\nu/N^2),
\]
independently. Put
\[
 Q=\frac dN\sum_j\|w_j\|_2^2+N\sum_jv_j^2,
 \quad P=\sqrt d\sum_j|v_j|\|w_j\|_2\le Q/2,
 \quad\omega=N/\nu.
\]
For $|Y_i|\le1$, $N\ge n$, $0<\nu\le1$, and $0<\lambda\le1/2$,
form the actual posterior proportional to
$e^{-\omega\widehat R(z)/(2\lambda)}d\mu$.
If $J_*$ denotes the minimum of empirical squared loss plus
$2\lambda$ times implementation mass, over all finite normalized-ridge
implementations, then
\[
 0\le\mathbb E_\Pi[\widehat R(z)+2\lambda P]-J_*
 \le C(d+2)\lambda\nu
          \log\frac{e(d+1)}{\lambda\nu}. \tag{GD1}
\]
The unique convex fitted training vector $b_*$ also satisfies
\[
 \mathbb E_\Pi\frac1n\|(z(X_i))_{i=1}^n-b_*\|_2^2
 \le C(d+2)\lambda\nu
          \log\frac{e(d+1)}{\lambda\nu}, \tag{GD2}
\]
and so does the training vector of the ordinary posterior mean.
\end{proposition}
\begin{proof}
Normalize a ridge as
$\phi_\theta(x)=d^{-1/2}(\theta\cdot\bar x)_+$ for
$\theta\in\mathbb S^d$. These features are bounded by $\sqrt2$,
and the signed amplitude is $c=\sqrt d\,v\|w\|_2$.
The compact signed-feature linear program again gives an optimum
$f_*$ with at most $n$ atoms and
$\widehat R(f_*)\le1$, $P_*\le(2\lambda)^{-1}$.
Balance an atom of amplitude $c$ with
\[
 \bar w=\sqrt{N|c|/d}\,\theta,
 \qquad \bar v=\operatorname{sign}(c)\sqrt{|c|/N}.
\]
The total center energy is $Q(\bar w)=2P_*$.

The parameter dimension is $D=(d+2)N$. A Gaussian comparison tube
with native covariance multiplied by $\epsilon^2$ has
\[
 \operatorname{KL}(\rho\Vert\mu)
 \le\omega P_*+D\log(1/\epsilon).
\]
Take $\epsilon=(\lambda\nu/(d+1))^2$.
The ReLU Lipschitz property, $\|\bar x\|_2^2\le d+1$, balance,
and independence of the incoming and readout perturbations give
uniform root-mean-square error at most
$C\epsilon\sqrt{\nu P_*}+C\epsilon^2\nu$.
The empirical loss increase is therefore at most $C\lambda\nu$.
The Gibbs variational principle gives
\[
 F\le J_*+C(d+2)\lambda\nu
                   \log\frac{e(d+1)}{\lambda\nu}.
\]

For $0<t<1$ the exact Gaussian moment is
\[
 \log\mathbb E_\mu e^{t\omega Q/2}
 =\frac D2\log\frac1{1-t}.
\]
The entropy inequality and $Q\ge2P$ imply
\[
 F\ge\mathbb E_\Pi\widehat R+2\lambda t\mathbb E_\Pi P
      -(d+2)\lambda\nu\log\frac1{1-t}.
\]
Gaussian radial integration by parts gives
\[
 \mathbb E_\Pi Q
 =(d+2)\nu-\frac2\lambda\mathbb E_\Pi P_n[z(z-Y)]
 \le(d+2)\nu+\frac1{2\lambda}.
\]
Use $1-t=\lambda\nu$ and $P\le Q/2$ to bound the lost coercivity
by $C(d+2)\lambda\nu$. This proves (GD1).
Strong convexity in the training vector proves (GD2) exactly as in
(PR4), and Jensen's inequality handles the posterior mean.
\end{proof}
